% Primary-source metadata and theorem locators checked on 7 October 2026.
% Self-contained bibliography: no BibTeX or Biber invocation is required.
\begin{thebibliography}{99}
\addcontentsline{toc}{section}{References}

\bibitem{barnatan}
Dror Bar-Natan,
\emph{Fast Khovanov homology computations},
Journal of Knot Theory and Its Ramifications
\textbf{16} (2007), no.~3, 243--255.
\href{https://arxiv.org/abs/math/0606318v1}{arXiv:math/0606318v1}.
\href{https://www.math.utoronto.ca/~drorbn/papers/FastKh/FastKh.pdf}{Author's version of 20 May 2007}.
Section~3 gives the local dotted relations; Lemmas~4.1--4.2 give
delooping and abstract Gaussian elimination.

\bibitem{barnatancobordisms}
Dror Bar-Natan,
\emph{Khovanov's homology for tangles and cobordisms},
Geometry \& Topology \textbf{9} (2005), 1443--1499.
\href{https://doi.org/10.2140/gt.2005.9.1443}{doi:10.2140/gt.2005.9.1443}.
\href{https://arxiv.org/abs/math/0410495v2}{arXiv:math/0410495v2}.
This supplies the formal tangle/cobordism framework underlying the local
algorithm; see Section~11.2 for the dotted specialization.

\bibitem{bhkk}
Andreas Bj\"orklund, Thore Husfeldt, Petteri Kaski, and Mikko Koivisto,
\emph{Fourier meets M\"obius: fast subset convolution},
in Proceedings of the 39th Annual ACM Symposium on Theory of Computing
(STOC 2007), ACM, 2007, 67--74.
\href{https://doi.org/10.1145/1250790.1250801}{doi:10.1145/1250790.1250801}.
\href{https://arxiv.org/abs/cs/0611101v1}{arXiv:cs/0611101v1}.
The algorithm computes subset convolution in $O(r^2 2^r)$ ring operations.

\bibitem{blair}
Ryan Blair, Alexandra Kjuchukova, Roman Velazquez, and Paul Villanueva,
\emph{Wirtinger systems of generators of knot groups},
Communications in Analysis and Geometry
\textbf{28} (2020), no.~2, 243--262.
\href{https://doi.org/10.4310/CAG.2020.v28.n2.a2}{doi:10.4310/CAG.2020.v28.n2.a2}.
\href{https://arxiv.org/abs/1705.03108v1}{arXiv:1705.03108v1}.
The introduction defines seed strands and diagrammatic propagation;
Theorem~1.3 identifies the minimum over diagrams with bridge number.

\bibitem{clark}
W.~Edwin Clark, Mohamed Elhamdadi, Masahico Saito, and Timothy Yeatman,
\emph{Quandle colorings of knots and applications},
Journal of Knot Theory and Its Ramifications
\textbf{23} (2014), no.~6, article 1450035, 29~pp.
\href{https://doi.org/10.1142/S0218216514500357}{doi:10.1142/S0218216514500357}.
\href{https://arxiv.org/abs/1312.3307v2}{arXiv:1312.3307v2}.
Appendix~D describes the normalization fixing one color in a connected
quandle braid-coloring search.

\bibitem{fish}
Andrew Fish, Alexei Lisitsa, and David Stanovsk\'y,
\emph{A combinatorial approach to knot recognition},
preprint, 25 May 2015.
\href{https://arxiv.org/abs/1505.06595v1}{arXiv:1505.06595v1}.
Theorems~2--3 discuss finite-quandle detection and the relation between
finite affine colorings and the Alexander polynomial; the paper also
describes SAT-based coloring search.

\bibitem{kelomaki}
Tuomas Kelom\"aki and Dirk Sch\"utz,
\emph{On computational complexity of Khovanov homology},
preprint, 5 January 2026, 30~pp.
\href{https://arxiv.org/abs/2601.02119v1}{arXiv:2601.02119v1}.
Theorems~1.2--1.4 distinguish exponential scanning examples, the specialized
three-braid algorithm, and extremal homological bands. Proposition~5.1 and
Lemma~5.5 give the supplied-nice-scan bounds used in the parameter discussion.

\bibitem{kronheimer}
Peter B.~Kronheimer and Tomasz S.~Mrowka,
\emph{Khovanov homology is an unknot-detector},
Publications Math\'ematiques de l'IH\'ES
\textbf{113} (2011), 97--208.
\href{https://doi.org/10.1007/s10240-010-0030-y}{doi:10.1007/s10240-010-0030-y}.
\href{https://arxiv.org/abs/1005.4346v1}{arXiv:1005.4346v1}.
Corollary~1.3, Proposition~1.4, and the discussion in the introduction give
reduced rational rank greater than one for every nontrivial knot. The
universal coefficient theorem then gives the required lower bound over
$\mathbb F_2$.

\bibitem{kuperbergsamperton}
Greg Kuperberg and Eric Samperton,
\emph{Coloring invariants of knots and links are often intractable},
Algebraic \& Geometric Topology
\textbf{21} (2021), 1479--1510.
\href{https://doi.org/10.2140/agt.2021.21.1479}{doi:10.2140/agt.2021.21.1479}.
\href{https://arxiv.org/abs/1907.05981v1}{arXiv:1907.05981v1}.
Theorem~1.1 gives the counting hardness result; Corollary~1.2 gives
NP-completeness of nonconstant coloring for a fixed nontrivial conjugacy
class in a fixed finite nonabelian simple group.

\bibitem{lackenbytalk}
Marc Lackenby,
\emph{Unknot recognition in quasi-polynomial time},
Oxford topology seminar handout, March 2021, 34~pp.
\href{https://people.maths.ox.ac.uk/lackenby/quasipolynomial-talk-oxford-compressed.pdf}{Author's seminar handout}.
The announced $2^{O((\log n)^3)}$ bound appears on slide~7;
the proposed speed-up ingredients are summarized from slide~25 onward.
This citation records the announcement and does not turn the slides into
a complete published complexity proof.

\bibitem{lackenby2026}
Marc Lackenby,
\emph{Incompressible surfaces, hierarchies and unknot recognition},
preprint, 25 July 2026, 54~pp.
\href{https://arxiv.org/abs/2607.23350v1}{arXiv:2607.23350v1}.
\href{https://arxiv.org/html/2607.23350v1}{Version-pinned full text}.
Proposition~9.1 gives the $L(g+1)^L$ step-visit bound. Theorem~14.1
concerns polynomially encoded and verifiable essential hierarchies.
The introduction leaves additional running-time improvements to future work.

\bibitem{lackenbycurves}
Marc Lackenby,
\emph{Some fast algorithms for curves in surfaces},
Discrete \& Computational Geometry, published online 21 April 2026.
\href{https://doi.org/10.1007/s00454-026-00845-7}{doi:10.1007/s00454-026-00845-7}.
\href{https://arxiv.org/abs/2401.16056}{arXiv:2401.16056}.
\href{https://people.maths.ox.ac.uk/lackenby/AlgorithmsSurfaces13jan26.pdf}{Author's version of 13 January 2026}.
The compressed normal-curve primitives are relevant to the geometric
research programme, without supplying a complete hierarchy-search bound.

\bibitem{openai}
OpenAI,
\emph{math: mathematical manuscripts and catalogue},
repository snapshot at commit
\texttt{adc7f1241b42e322a6451854ab7e4b4c146bf78a},
accessed 7 October 2026.
\href{https://github.com/openai/math/tree/adc7f1241b42e322a6451854ab7e4b4c146bf78a}{Pinned repository}.
Only the README and catalogue were used for the applicability audit;
no result in this article relies on an unchecked manuscript from the collection.

\bibitem{proveit}
Vladimir Reshetnikov,
\emph{ProveIt: Topology/UnknotRecognition},
software and research repository, source snapshot at commit
\texttt{ca61a1a5f967999d778d6efd7c882e77a2eb379d},
accessed 7 October 2026.
\href{https://github.com/VladimirReshetnikov/ProveIt/tree/ca61a1a5f967999d778d6efd7c882e77a2eb379d/Topology/UnknotRecognition}{Pinned project source}.
The present integration patch targets this snapshot's \texttt{fast/} subtree.

\bibitem{report07}
ProveIt research continuation,
\emph{Unknot recognition by structural certificates and compressed chain
complexes: exact improvements and a finite-type route towards
quasi-polynomial time},
research report~07, 7 October 2026.
\href{https://github.com/VladimirReshetnikov/ProveIt/tree/ca61a1a5f967999d778d6efd7c882e77a2eb379d/Topology/UnknotRecognition/reports/07}{Report and accompanying implementation in the pinned repository}.
This is a project research artifact, not a journal publication. Its own
implementation baseline is identified as commit
\texttt{4e6fe879e5238ec0b134b6d33fdca0b8c7c1c711}.

\bibitem{shumakovitch}
Alexander N.~Shumakovitch,
\emph{Torsion of Khovanov homology},
Fundamenta Mathematicae \textbf{225} (2014), 343--364.
\href{https://doi.org/10.4064/fm225-1-16}{doi:10.4064/fm225-1-16}.
\href{https://arxiv.org/abs/math/0405474v2}{arXiv:math/0405474v2},
revised 19 June 2018 to match the published version.
Corollary~3.2.C gives the direct-sum relation between unreduced and reduced
Khovanov homology over $\mathbb F_2$, with quantum shifts $\pm1$.

\end{thebibliography}
